\scanpage{62}
% upright top left, lines 1–8
% letterhead
U N I V E R S I T E\\
DES SCIENCES ET TECHNIQUES\\
DU LANGUEDOC\\
\texttt{------------}\\
Service de la Scolarité\\
MR / GJ\\
Diffusion : Liste A\\
N° : 4151

% upright top right, line 1
\unsure{-} Montpellier, le 1er Juillet 1978

% upright upper centre, lines 9–10
% printed form
\uline{AVIS DE SOUTENANCE D'UNE THESE DE DOCTORAT DE 3ème CYCLE}

\uline{(PHYSIQUE DU SOLIDE\quad - APPLICATION A L'ELECTRONIQUE\qquad )}

% upright, field 1
\uline{Nom et prénom de la candidate} : MERLET née BERNARD Catherine

% upright, field 2
\uline{Titre de la thèse} : Le phénomène de gel magnétique dans InSb : étude en fonction du dopage.

% upright, field 3
\uline{Date, heure et lieu de soutenance prévus} :\\
12 JUILLET 1978 à 15 Heures,\\
Salle de Séminaire du C.E.E.S.- U.S.T.L.

% upright, field 4
\uline{Nom du Directeur de thèse} :\\
Monsieur le Professeur ROBERT.

% upright, field 5
\uline{Désignation du laboratoire où la thèse est préparée} :\\
Centre d'Etudes d'Electronique des Solides, MONTPELLIER.

% upright, field 6
\uline{Analyse de l'ouvrage} :\\
L'étude du phénomène de gel magnétique a été faite sur des échantillons d'InSb différemment dopés ($10^{14}$à $10^{16}\mathrm{cm}^{-3}$). Ce phénomène se traduit par une variation de l'énergie de liaison en fonction du dopage. Pour les échantillons faiblement dopés on a mis en évidence la localisation d'une partie des porteurs de la bande de conduction à faible température.

% upright, field 7
\uline{Durée approximative des recherches ayant abouti à la thèse} : Deux ans

% upright, field 8
\uline{Composition du jury proposé} :\\
MM. :\\
PISTOULET, Président\\
ROBERT\\
RAYMOND\\
LEROUX HUGON\\
THUILLIER

% upright lower right, closing
Le Secrétaire Général,\\

% upright lower right, signature
\handwritten{\unsure{E Siau}}\\

% upright lower right, typed name
\uline{E. SIAU}

% foot, left, in supplied orientation
% archive sheet number: 61

